\section{Linear-momentum density}

The electromagnetic linear-momentum density is related directly to the
Poynting vector by
\begin{align}
  {\bf g}(t,{\bf x})=\frac1{c^2}{\bf S}(t,{\bf x})
  =\epsilon_0{\bf E}(t,{\bf x})\times{\bf B}(t,{\bf x}).
  \label{eq:linear-momentum-density}
\end{align}
Therefore no separate derivation is needed.  Dividing
Eq.~\eqref{eq:Poynting-vector-Faraday} by $c^2$ gives
\begin{align}
  g^i(t,{\bf x})
  =-\frac1{\mu_0c}\sum_{j=1}^{3}
  F^{j0}(t,{\bf x})F^{ij}(t,{\bf x}).
  \label{eq:linear-momentum-density-Faraday}
\end{align}
Similarly, Eq.~\eqref{eq:instantaneous-Poynting-vector-Fourier-F} gives the
instantaneous momentum density in terms of the Fourier-transformed Faraday
tensor,
\begin{align}
  g^i(t,{\bf x})
  =-\frac1{\mu_0c}
  \int_{-\infty}^{\infty}d\omega
  \int_{-\infty}^{\infty}d\omega'\,
  e^{-i(\omega+\omega')t}
  \sum_{j=1}^{3}
  F^{j0}(\omega,{\bf x})F^{ij}(\omega',{\bf x}).
  \label{eq:instantaneous-linear-momentum-density-Fourier-F}
\end{align}

The time-integrated, two-sided spectral momentum density is simply
\begin{align}
  &\int_{-\infty}^{\infty}dt\,g^i(t,{\bf x})
  =\int_{-\infty}^{\infty}d\omega\,
  \frac{d\mathcal G^i}{d\omega}(\omega,{\bf x}),\\
  &\frac{d\mathcal G^i}{d\omega}(\omega,{\bf x})
  =\frac1{c^2}\frac{d\mathcal S^i}{d\omega}(\omega,{\bf x})
  =-\frac{2\pi}{\mu_0c}\operatorname{Re}
  \left[
    \sum_{j=1}^{3}F^{j0}(\omega,{\bf x})
    F^{ij}(\omega,{\bf x})^*
  \right].
  \label{eq:two-sided-spectral-linear-momentum-density-F}
\end{align}
\begin{highlightbox}[title={One-sided spectral linear-momentum density}]
  For the one-sided spectrum,
  \begin{align}
    \frac{d\mathcal G^i_+}{d\omega}(\omega,{\bf x})
    =\frac1{c^2}\frac{d\mathcal S^i_+}{d\omega}(\omega,{\bf x})
    =-\frac{4\pi}{\mu_0c}\operatorname{Re}
    \left[
      \sum_{j=1}^{3}F^{j0}(\omega,{\bf x})
      F^{ij}(\omega,{\bf x})^*
    \right],\qquad \omega>0.
    \label{eq:one-sided-spectral-linear-momentum-density-F}
  \end{align}

\end{highlightbox}
